
Semantic entropy (Kuhn et al., 2023) computes entropy over semantically equivalent model outputs by clustering answers with similar meanings before calculating uncertainty. The method was validated on natural language generation tasks and showed improved hallucination detection compared to token-level entropy. The original work did not include an ablation isolating the clustering contribution from the computational cost of multiple samples.

Self-consistency (Wang et al., 2022) samples multiple outputs and measures agreement, using the most common answer as the final prediction. Originally developed for chain-of-thought reasoning on mathematical and commonsense tasks, the method showed substantial accuracy improvements. The work focused on reasoning tasks rather than factual error detection and did not analyze orthogonality with other uncertainty methods.

Verbalized confidence (Kadavath et al., 2022) prompts models to self-report confidence levels. Their work demonstrated that language models can provide reasonably calibrated confidence estimates when explicitly asked. The analysis characterized calibration properties but did not compare verbalized confidence to output-based methods or identify conditions determining calibration quality across benchmarks.

Token probability methods use metrics like entropy or variance over token-level probability distributions. These methods require access to model logits but are computationally efficient, requiring only a single forward pass.

TruthfulQA (Lin et al., 2021) tests whether models generate truthful answers to questions designed to elicit common misconceptions. NaturalQuestions (Kwiatkowski et al., 2019) provides factual questions from Google search queries with Wikipedia-based answers, including an unanswerable subset.

Our work differs from prior uncertainty estimation research by prioritizing mechanistic understanding through ablation studies and correlation analysis, by treating null results as contributions to knowledge, and by grounding claims in controlled experiments rather than empirical comparisons where multiple factors vary simultaneously.
